\chapter*{Lời Mở Đầu}
\addcontentsline{toc}{chapter}{Lời Mở Đầu}
\markboth{Lời Mở Đầu}{Lời Mở Đầu}

\begin{truyen}{Lời Mở Đầu}{Opening Words}
\chuhoa{C}{ó những điều học sinh không cần đợi lớn mới biết là đúng.}
\textit{(There are things students do not need to wait until adulthood to know are true.)}

Nếu hiểu sớm, các em sẽ đỡ tự làm khó mình hơn trong việc học, trong cách nghĩ về bản thân, trong chuyện bạn bè và trong việc đi về phía tương lai.
\textit{(If understood early, students will make life less difficult for themselves in study, in self-image, in friendship, and in moving toward the future.)}

Những điều trong quyển này không hẳn mới.
\textit{(The things in this book are not entirely new.)}

Nhưng nhiều điều tuy quen vẫn cần được nói lại cho rõ, cho gần, và cho đủ sớm.
\textit{(But many familiar truths still need to be said again clearly, closely, and early enough.)}

Quyển sách này vì thế không cố làm điều lớn lao. Nó chỉ gom lại những nền tảng quan trọng mà một học sinh hiểu càng sớm thì càng đỡ vấp về sau.
\textit{(This book does not aim to be grand. It simply gathers important foundations that, if understood early, can save students from many later stumbles.)}

Nếu em đang đi học, mong em đọc từng bài như một lời nhắc ngắn nhưng thật.
\textit{(If you are a student, I hope you read each piece as a short but honest reminder.)}

Nếu em đang dạy học, mong quyển này giúp mình nói lại những điều căn bản bằng một giọng bình tĩnh hơn.
\textit{(If you are a teacher, I hope this helps us say foundational things again in a calmer voice.)}
\end{truyen}

\begin{baihoc}
Có những điều càng hiểu sớm, con người càng đỡ đi vòng hơn trong việc học và lớn lên.
\textit{(Some things, the earlier we understand them, the fewer unnecessary detours we take in learning and growing.)}
\end{baihoc}

\begin{ghinhoanh}
\textbf{Short English to remember:}

Early understanding can prevent later confusion.

\end{ghinhoanh}

\begin{tuhocnhanh}
1. Vì sao có những điều cần hiểu sớm? \textit{Why are some truths worth understanding early?}

2. Quyển này muốn giúp học sinh bớt vấp ở đâu? \textit{Where does this book hope to reduce later stumbles?}

3. Điều gì em ước mình hiểu sớm hơn? \textit{What do you wish you had understood earlier?}
\end{tuhocnhanh}

\ngancach